\chapter*{Introduction générale}
\addcontentsline{toc}{chapter}{Introduction générale}
\markboth{Introduction générale}{}

La valorisation des produits dérivés est un enjeu central de l'ingénierie financière et de l'assurance : banques et assureurs doivent attribuer un prix à des contrats dont la valeur dépend de l'évolution future d'un actif, le \emph{sous-jacent}. Depuis les travaux de Black, Scholes et Merton (1973), la modélisation stochastique est le cadre de référence.

Ce projet du module \emph{Mathématiques et Assurance} étudie de façon méthodique la \textbf{simulation de Monte-Carlo}. Il procède en deux temps. L'\textbf{option d'achat (call) européenne} sert d'abord de \emph{banc d'essai} idéal : son prix exact $\CBS$ est donné par la formule de Black-Scholes, ce qui permet de mesurer exactement l'erreur de la simulation, sa vitesse de convergence et l'efficacité des techniques de réduction de variance. La méthode, ainsi validée, est ensuite appliquée à une option dépendante du chemin, l'\textbf{option asiatique} à moyenne arithmétique, qui n'a pas de formule fermée : c'est là que Monte-Carlo montre sa vraie valeur pratique.

\paragraph{Démarche.} Le travail suit les quatre étapes de l'énoncé :
\begin{enumerate}[nosep]
  \item modèle de Black-Scholes, résolution log-normale et formule analytique $\CBS$ ;
  \item discrétisation du mouvement brownien géométrique (schéma exact et schéma d'Euler), Monte-Carlo standard (la « méthode de base »), avec des tirages gaussiens générés par la méthode de Box-Muller, et intervalle de confiance à 95\,\% pour $10^3\le N\le 10^6$ ;
  \item optimisation par réduction de variance (variables antithétiques et variable de contrôle), comparaison de l'erreur quadratique moyenne, du gain de variance et du temps de calcul ;
  \item option asiatique : simulation de la grille temporelle complète, prix, sensibilité à la fréquence $M$ et au nombre de trajectoires $N$.
\end{enumerate}

\paragraph{Organisation du rapport.}
Le \textbf{chapitre 1} définit les options et les classe selon deux critères, en commençant par la forme du payoff.
Le \textbf{chapitre 2} présente d'abord les \emph{modèles} du sous-jacent (dont Black-Scholes, démontré en détail), puis les \emph{méthodes} de valorisation, chacune expliquée et illustrée sur le call.
Le \textbf{chapitre 3} réunit la discrétisation, la simulation de Monte-Carlo et la réduction de variance, avec les démonstrations (dont celle du coefficient optimal $c^*$ de la variable de contrôle).
Le \textbf{chapitre 4} applique la méthode à l'option asiatique.
Le \textbf{chapitre 5} explique le code partie par partie.
L'\textbf{annexe A} indique où chaque remarque de l'encadrant a été traitée.

\paragraph{Outils.} Le code est un notebook Jupyter Python (\texttt{numpy}, \texttt{scipy} et \texttt{matplotlib}) entièrement reproductible : chaque partie a son générateur aléatoire à graine fixée, et tous les nombres et figures de ce rapport sont produits par ce notebook. Paramètres de référence :
\[
S_0 = 100,\qquad K = 100,\qquad T = 1 \text{ an},\qquad r = 5\,\%,\qquad \sigma = 20\,\%.
\]
